\documentclass[10pt,a4paper]{amsart}
\usepackage[margin=19mm]{geometry}
\usepackage{amsmath,amssymb,mathtools}
\usepackage[scr=rsfs,cal=euler]{mathalfa}
\usepackage{xfrac}
\usepackage[unicode,hidelinks,bookmarks=false]{hyperref}
\newcommand{\ie}{i.e.\ }
\newcommand{\cf}{cf.\ }
\newcommand{\resp}{respectively\ }
\newcommand{\Subsection}{\subsection}
\providecommand{\nobreakdash}{\nobreak\hbox{-}}
\setlength{\emergencystretch}{2em}
\multlinegap0pt
\usepackage{xcolor}
\usepackage{caption}
\usepackage{amssymb}
\usepackage[all]{xy}
\usepackage[normalem]{ulem}
\usepackage{tcolorbox}
\usepackage{xfrac}
\newtheorem{proposition}{Proposition}
\newtheorem{lemma}{Lemma}
\newtheorem{theorem}{Theorem}
\usepackage{cleveref}
\crefname{proposition}{Proposition}{Propositions}
\crefname{lemma}{Lemma}{Lemmas}
\crefname{theorem}{Theorem}{Theorems}
\DeclareMathOperator{\Ext}{Ext}
\newcommand{\cB}{{\mathcal B}}
\newcommand{\cE}{{\mathcal E}}
\newcommand{\mF}{{\mathbb F}}
\newcommand{\mM}{{\mathbb M}}
\newcommand{\mR}{{\mathbb R}}
\title{A spectrum-level splitting of the $ku_\mathbb{R}$-cooperations algebra}
\author{Guchuan Li, Sarah Petersen, Elizabeth Tatum}
\date{Source: arXiv:2503.17149v3 (CC BY 4.0)}
\begin{document}
\maketitle

\noindent\emph{Source and licence.} A spectrum-level splitting of the $ku_\mathbb{R}$-cooperations algebra. Version arXiv:2503.17149v3 (CC BY 4.0), distributed by arXiv under the Creative Commons Attribution 4.0 International licence, http://creativecommons.org/licenses/by/4.0/, as recorded in the arXiv licence metadata for this exact version and shown on the abstract page https://arxiv.org/abs/2503.17149v3. Full licence notice: \texttt{licenses/arxiv-2503.17149-CC-BY-4.0.txt}.\par\smallskip

\noindent\textit{Editorial scope.} Counted unit: the article's proposition prop:splitInto4 (source0.tex lines 1343--1352). The article's complete original proof of it is delivered with it (source0.tex lines 1779--1926, one \texttt{proof} environment of 148 lines, which the article places away from the statement). No mathematics is written, repaired or re-derived: the statement and the proof are the article's own lines, in the article's own notation, and no retained line is substituted or re-flowed.  Retained uncounted as the source-local context the proof needs: lines 462--464; lines 513--528; lines 543--549; lines 1386--1413; lines 1531--1535; lines 1596--1636; lines 1713--1717; lines 1750--1757.  One declared adaptation: exactly 4 ancillary strings inside the retained spans are omitted (image-inclusion commands and, where the source repeats a label, the repeated label definitions) and no other line differs from the source; the figure environments, with their placement, captions and labels, are kept, and the package records the omitted strings in fidelity receipts bound to the affected bodies.  No retained line contains a citation, so the wrapper carries no bibliography and no bibliography entry.  The retained lines contain the article's own diagram code, which the wrapper typesets with the standard tikz packages; no image file is delivered and the package declares no asset dependency.  Editorial formatting only, in addition: an AMS \texttt{amsart} preamble that restates the article's own declarations and macros used by the retained lines, namely the environments \texttt{proposition}, \texttt{lemma}, \texttt{theorem} on separate counters, together with the standard packages those lines need.  The retained environments print in this document as Proposition 1, Lemma 1, Theorem 1 in order of appearance, whereas the article prints the same items with the numbering of its own full document.  The complete native source of the article as distributed in the arXiv e-print for this exact version is bundled unchanged as \texttt{sources/175158/source0.tex}.
\begin{equation}\label{nonequiv:ass}
     \Ext_{E(1)_{*}}^{s,f} \left(H_{*} \Sigma^{2k} B_{0}(k), H_{*}ku \right) \Longrightarrow [ku \wedge \Sigma^{2k}B_{0}(k), ku \wedge ku]^{ku},
\end{equation}

\begin{figure}
\centering
\begin{minipage}{.5\textwidth}
\centering

 \captionof{figure}{\\$\Ext_{E(1)_*} (\mF_2, L(3))$}
  \label{fig:clExtFpL3}
\end{minipage}%
\begin{minipage}{.5\textwidth}
\begin{center}

 \end{center}
  \captionof{figure}{\\ $\Ext_{E(1)_*} (L(3), \mF_2)$}
  \label{fig:clExtL3Fp}
\end{minipage}
\end{figure}

\begin{proposition}\label{nonequiv: collapse}
    For all $k$, the Adams spectral sequence of (\ref{nonequiv:ass}), 
    \[
    \Ext_{E(1)_{*}}^{s,f} \left(H_{*} \Sigma^{2k} B_{0}(k), H_{*}ku \right) \Longrightarrow [ku \wedge \Sigma^{2k}B_{0}(k), ku \wedge ku]^{ku},
    \]
    collapses.
\end{proposition}

where the stem $s,$ filtration $f,$ and motivic weight $w$ of the generators are given in \cref{tab:ExtM2M2gens}. The negative cone summand, $\Ext_{\cE_\star^{-} (1)} (\mM_2, \mM_2):= \Ext_{\cE_{\star}^{\mR} (1)} (NC, \mM_2^{\mR})$, is a module over $\Ext_{\cE_\star^{+} (1)} (\mM_2, \mM_2)$. We list the generators of $\Ext_{\cE_\star^{-} (1)} (\mM_2, \mM_2)$ below the horizontal line in \cref{tab:ExtM2M2gens}.  There are relations of the form $v_{1}\frac{\gamma}{\tau^{4i + 3}} = v_{0}\frac{\gamma}{\rho^{2}\tau^{4i + 2}}$ for all $i$. The other relations will be clear from the names of the generators (for example, $\tau^{2}v_{0}\frac{\gamma}{\tau^{3}} = \frac{\gamma}{\tau} v_{0}$).
\begin{table}[ht]
    \centering
    \begin{tabular}{l l l}
          $(s,\, f,\, w)$ & Element & \\
            \hline
             $\,$ & $\,$ \\
             $(-1, 0, -1)$ & $\rho$ \\
             $(0, 1, 0)$ & $v_0$ \\
            $(2, 1, 1)$ & $v_1$ \\
            $(0, 1, -2)$ & $\tau^2 v_0$ \\
            $(0, 0, -4)$ & $\tau^4$ \\
           $\,$ & $\,$ \\
           \hline
           \, & \, \\
           $(0,0, 4j+2)$ & $\frac{\gamma}
           {\tau^{4j + 1}}$ & $0 \le j$\\
           $\left(2,0,4(j+1) + 1\right)$ & $\frac{\gamma}
           {\tau^{4j + 2}\rho^{2}}$ & $0 \le j$ \\
           $(0,0,4j + 4)$ & $\frac{\gamma}
           {\tau^{4j + 3}}$ & $0 \le j$ \\
           $(i,0,  4j  + i + 1)$ & $\frac{\gamma}
           {\tau^{4j}\rho^{i}}$ & $0 < j, 0 \le i$ \\
           \label{ExtM2gens}
    \end{tabular}
    \caption{Generators for $\Ext_{\cE (1)_\star} (\mM_2, \mM_2)$. The generators above the horizontal line are the generators of the ring $\Ext_{\cE^{+}_\star (1)} (\mM_2, \mM_2)$. The elements below the horizontal line generate $\Ext_{\cE^{-}_\star (1)} (\mM_2, \mM_2)$ as an $\Ext_{\cE^{+}_\star (1)} (\mM_2, \mM_2)$-module.}
    \label{tab:ExtM2M2gens}
\end{table}

\begin{figure}[ht]

\caption{$E_1$-page computing $\Ext_{\cE(1)_\star}(L(1), L(2))$}
\label{fig:E1ExtL1L2}
\end{figure}

\begin{lemma}\label{lem:k>m}
The positive cone $\Ext_{\cE^{+}_\star(1)}^{*,*,*}(L(k), L(m))$ when $k > m$ consists of: 
\begin{enumerate}
\item a triangle formation consisting of \label{classes:positive triangle}
\[
\Ext_{\cE^{+}_{\star}(1)}\left(\mM_{2}, \mM_{2}\right) \left\{y_{0}, \ldots, y_{m-k-1} \right\}
\]
with relations $v_{1}y_{i} = v_{0}y_{i+1},$ $v_{0}y_{0} = 0,$ and $v_{1}y_{m-k-1} = 0$ and 
\[
\Ext_{\cE^{+}_{\star}(1)}\left(\mM_{2}, \mM_{2}\right) \left\{ \tau^{2}y_{0}, \ldots, \tau^{2}y_{m-k-1} \right\}
\]
with relations $v_{1} \tau^{2}y_{i} = v_{0}\tau^{2}y_{i+1},$ $v_{0}^{2}\tau^{2}y_{0} = 0,$ and $v_{1}\tau^2 y_{m-k-1} = 0$ where \newline $|y_{i}| = \big(-2(k-m - i)-1, 0, -(k-m - i)\big),$ 

     \item infinite $\rho$-towers generated in odd stem

      \item a copy of $v_{1}\Ext_{\cE^{+}_{\star}(1)}(\mM_{2}, \mM_{2})$, with generator denoted $x$ and $|x| = (0, k-m, 0),$\label{component 3 of pos cone}

    \item $\rho$-pairs: \[\mF_{2}[\rho, v_{1}] \left\{ v_1 b \, | \, v_{1}^{m-k-1}b = \rho x, \, \rho^{2}b = 0\right\}, \] where $|b| = (2(m-k)-1, \, 0, \, m-k-1)$. 
    \end{enumerate}

    The negative cone $\Ext_{\cE^{-}_\star(1)}^{*,*,*}(L(k), L(m))$ when $k > m$ consists of:


\begin{enumerate} 
\item a triangle formation consisting of \label{classes:negative triangle}
\[ 
\Ext_{\cE^{+}_{\star}(1)}\left(\mM_{2}, \mM_{2}\right) \left\{ \frac{\gamma}{\tau^{4j + 3}}y_0, \cdots, \frac{\gamma}{\tau^{4j + 3}}y_{m-k-1} \right\} 
\]
with relations $v_{1}\frac{\gamma}{\tau^{4j + 3}}y_{i} = v_{0}\frac{\gamma}{\tau^{4j + 3}}y_{i+1},$ $v_{1}\frac{\gamma}{\tau^{4j + 3}}y_{m-k-1} = 0,$ and $v_{0}\frac{\gamma}{\tau^{4j + 3}}y_{0} = 0,$ and 
\[
\Ext_{\cE^{+}_{\star}(1)}\left(\mM_{2}, \mM_{2}\right) \left\{ \frac{\gamma}{\tau^{4j + 1}}y_0, \cdots, \frac{\gamma}{\tau^{4j + 1}}y_{m-k-1} \right\}
\]
with relations $v_{1}\frac{\gamma}{\tau^{4j + 1}} = v_{0}\frac{\gamma}{\tau^{4j + 1}}y_{i+1},$ $v_{1}\frac{\gamma}{\tau^{4j + 1}}y_{m-k-1} = 0,$ and $v_{0}\frac{\gamma}{\tau^{4j + 1}}y_{0} = 0$ where $|y_{i}| = \big(-2(k-m - i)-1, 0, -(k-m - i)\big),$ 
\item infinite $\rho$-divisible towers in filtration $f = 0$, as well as copies of $\mF_{2}[\tau^{2}]/\tau^{\infty}$ in odd stem,
\item a copy of the $\Ext_{\cE^{+}_{\star}(1)}(\mM_{2}, \mM_{2})$-submodule of $\Ext_{\cE^{-}_{\star}(1)}(\mM_{2}, \mM_{2})$ generated by \newline $\mF_{2}\frac{[\tau^{4}]}{\tau^{\infty}} \left\{ \frac{\gamma}{\rho^{2}\tau^{2}} \right\}$, with generator denoted by $x'$ and shifted so that $|x'| = (0, k-m-1, 4),$ 
\item $\rho$-pairs, with generator denoted by $c$ and $|c| = (2(m-k) +1, 0, m-k + 1)$ for $m-k > 1$
\[
\frac{\mF_2[\tau^4]}{\tau^\infty}[\rho, v_{1}]\{c | \rho^{2}c = 0,\ v_1^nc = x'\}.
\]
\end{enumerate}
\end{lemma}

\begin{figure}[ht]
    \centering

\caption{$\Ext_{E(1)_\star}(L(2), \mathbb{M}_2)$} \label{fig:ExtL2M2}
\end{figure}

\begin{theorem}\label{prop: spec seq collapse}
    For all $k \ge 0$, the Adams spectral sequence
   \begin{align*}
E_2^{s, f, w}  = Ext_{\mathcal{E}(1)_\star} \left(H_\star \Sigma^{\rho k} \cB_0(k), H_\star ku_\mR\right) 
      \implies \left[ku_\mR \wedge \Sigma^{\rho k} \cB_{0} (k), ku_\mR \wedge ku_\mR\right]^{ku_\mR}\notag
\end{align*}
collapses.
\end{theorem}

\begin{proposition} \label{prop:splitInto4}
   The Adams spectral sequence  
\[
 E_2^{s, t} = \Ext_{\mathcal{E}(1)_\star} \left(H_\star \Sigma^{\rho k} B_0(k), H_\star ku_\mR \right) \implies [ku_\mR \wedge \Sigma^{\rho k} B_{0} (k),ku_\mR \wedge ku_\mR]^{ku_\mR}
\]decomposes into a sum of four separate spectral sequences, listed below. All but the first is concentrated on the $(f=0)$-line.
\[ \Ext_{\cE(1)_{\star}}(H_{\star}\Sigma^{\rho k}C_{k}, H_{\star}C) \Longrightarrow [\Sigma^{\rho k}C_{k}, C], \]
\[ \Ext_{\cE(1)_{\star}}(H_{\star}\Sigma^{\rho k}C_{k}, H_{\star}V) \Longrightarrow [\Sigma^{\rho k}C_{k}, V], \]
\[ \Ext_{\cE(1)_{\star}}(H_{\star}\Sigma^{\rho k}V_{k}, H_{\star}C) \Longrightarrow [\Sigma^{\rho k}V_{k}, C], \]
\[ \Ext_{\cE(1)_{\star}}(H_{\star}\Sigma^{\rho k}V_{k}, H_{\star}V) \Longrightarrow [\Sigma^{\rho k}V_{k}, V] .\]
\end{proposition}

\begin{proof}[Proof of \cref{prop: spec seq collapse}]
In \cref{prop:splitInto4}, we showed that the Adams spectral sequence splits into a sum of four different Adams spectral sequences, and that all of the summands collapse at the $E_2$-page (and are concentrated on the $(f = 0)$-line) except  
\[ 
\Ext_{\mathcal{E}(1)_\star}^{s,f,w} (H_\star \Sigma^{\rho k}C_{k}, H_\star C) \Longrightarrow [\Sigma^{\rho k}C_{k}, C] .
\]
We also observed
\[ 
\Ext_{\mathcal{E}(1)_\star}^{s,f,w} (H_\star\Sigma^{\rho k} C_{k}, H_\star C) \cong  \bigoplus\limits_{m=0}^\infty \Sigma^{\rho (m-k)} \Ext_{\mathcal{E}(1)_\star} (L (\nu_2(k!)), L(\nu_2(m!))) .
\] 
To show the Adams spectral sequence collapses, we will begin by considering Adams differentials with source a $v_{1}$-torsion class.

\subsubsection{Differentials originating in $v_{1}$-torsion classes} The $v_{1}$-torsion classes in \[
\bigoplus\limits_{m=0}^\infty \Sigma^{\rho (m-k)} \Ext_{\mathcal{E}(1)_\star} (L (\nu_2(k!)), L(\nu_2(m!)))
\]
consist of: 
\begin{enumerate}
    \item $v_0$-towers in the triangle formation (\cref{lem:k>m} (1)),
    \item $\rho$-towers based in odd stem and filtration $f = 0$ (\cref{lem:k>m} (2)),
    \item $\rho$-divisible towers in filtration $f=0$ (\cref{lem:k>m} (2)). 
\end{enumerate} 

By $v_1$-linearity, the only potential targets of differentials with source a $v_1$-torsion class are other $v_1$-torsion classes. All $v_1$-torsion classes in filtration $f > 0$ are contained in the $v_0$-towers in the triangle formation and are $\rho$-torsion. Hence the $\rho$ and $\rho$-divisible towers cannot support differentials. Further, the $v_0$-towers in the triangle formation are all in odd stem, so they cannot support differentials as Adams differentials decrease stem by one. 

Before analyzing the possible differentials originating in $v_{1}$-torsion free classes, it will be helpful to list the generators of $Ext_{\cE(1)_{\star}}(L(k), L(m))$ as a module over the positive cone summand $\Ext_{\cE^{+}_\star}(1)((L(k), L(m))$. 
 
\subsubsection{Listing Generators.} Since differentials with source a $v_1$-torsion free class can have target a $v_1$-torsion class, we consider both $v_1$-torsion and $v_1$-torsion free classes. However, we omit $\rho$ and $\rho$-divisible towers from the remaining analysis since, as they are concentrated in filtration $f=0,$ they cannot be the target of any differentials.

\cref{tab:Extgens k<m pos} lists the generators of the positive cone summand of $\Ext_{\cE (1)_\star} (L(k), L(m))$ as a module over the positive cone summand $\Ext_{\cE^{+}_{\star}(1)}(1)((L(k), L(m))$ when $k \le m$. \cref{tab:Extgens k<m neg} lists the generators of the negative cone summand. It may be helpful to compare these tables with the chart for $\Ext_{\cE (1)_\star} (L(1), L(2))$ in \cref{fig:E1ExtL1L2}). 

\cref{tab:Extgens k>m pos} lists the generators of the positive cone summand of $\Ext_{\cE (1)_\star} (L(k), L(m))$ as a module over the positive cone summand $\Ext_{\cE^{-}_{\star}(1)}(1)((L(k), L(m))$ when $k > m.$ \cref{tab:Extgens k>m neg} lists the generators of the negative cone summand. In both tables, the $v_{1}$-torsion free generators are listed above the horizontal line, and $v_{1}$-torsion generators are below. It may also be helpful to compare these tables with the chart given for $\Ext_{\cE (1)_\star} (L(2), \mM_2)$ in \cref{fig:ExtL2M2}).

\begin{table}[h!] 
    \centering
    \begin{tabular}{l |l |l}
            Generators & $(s,\, f,\, w)$ & \\
            \hline
            $x_{i}$ & $(2i, 0, i)$ & $0 \le i \le m-k$ \\
            $\tau^{2}x_{i},$ & $(2i, 0, i-2)$ & $0 \le i < m-k$ \\
    \end{tabular}
    \caption{Generators of $\Ext_{\cE^{+}_{\star}(1)} (L(k), L(m))$ for $k \le m$, omitting $\rho$ and $\rho$-divisible towers}
    \label{tab:Extgens k<m pos}
\end{table}

\begin{table}[h!] 
    \centering
    \begin{tabular}{l |l |l}
            Generators & $(s,\, f,\, w)$ & \\
            \hline
            $\frac{\gamma}{\tau^{2j+1}}x_{i}$ & $(2i, 0, i + 2j + 2)$ & $0 \le i \le m-k$, $j \ge 0$ \\
            $\frac{\gamma}{\rho^{2}\tau^{2j}}x_{m-k}$ & $(2(m-k+1), 0, m-k + 2j + 3)$ & $j \ge1 $ \\
    \end{tabular}
    \caption{Generators of $\Ext_{\cE^{-}_{\star}(1)} (L(k), L(m))$ for $k \le m$, omitting $\rho$ and $\rho$-divisible towers}
    \label{tab:Extgens k<m neg}
\end{table}

\begin{table}[h!] 
    \centering
    \begin{tabular}{l |l |l}
            Generators & $(s,\, f,\, w)$ & \\
            \hline
            $x$ & $(0, k-m, 0)$ &  \\
            $b$ & $(2(m-k)-1, 1, m-k-1)$ & \\
            \hline
            $y_{i}$ &{$\left(-2(k-m-i) -1, 0, -(k-m-i)\right)$} &{$0 \le i < k-m$} \\
            $\tau^{2}y_{i}$ & {$\left(-2(k-m-i) -1, 0, -(k-m-i)   - 2\right)$} & {$0 \le i < k-m$}\\
    \end{tabular}
    \caption{Generators of $\Ext_{\cE^{+}_{\star}(1)} (L(k), L(m))$ for $k > m$, omitting $\rho$-towers and $\rho$-divisible towers} 
    \label{tab:Extgens k>m pos}
\end{table}

\begin{table}[h!] 
    \centering
    \begin{tabular}{l |l |l}
            Generators & $(s,\, f,\, w)$ & \\
            \hline
            $x'$  & $(0, k-m-1, 4)$ & {}\\
            $\frac{\gamma}{\tau^{4j}}c$
            & $(2(m-k)+1, 0, m-k + 4j+ 1)$ 
            & {} \\
            \hline
            $\frac{\gamma}{\tau^{2j + 1}}y_{i}$ &{$\left(-2(k-m-i) -1, 0, -(k-m-i) + 2j + 2\right)$} &{$0 \le i < k-m$} \\
            
    \end{tabular}
    \caption{Generators of $\Ext_{\cE^{-}_{\star}(1)} (L(k), L(m))$ for $k > m$, omitting $\rho$ and $\rho$-divisible towers}
    \label{tab:Extgens k>m neg}
\end{table}

\subsubsection{Differentials from $v_{1}$-torsion free classes to $v_{1}$-torsion free classes.}
Since the Adams differential decreases stem by $1$ but preserves motivic weight, it will follow immediately from the following proposition that there are no differentials from $v_{1}$-torsion free classes to $v_{1}$-torsion free classes. 

\begin{proposition}
    Let $x$ be a $v_{1}$-torsion free class in $E_{2}^{s,f,w}$. If $x$ is in stem $2n$, then the motivic weight of $x$ is congruent to $n$ mod $2$. If $x$ is in stem $2n-1$, then the motivic weight of $x$ is congruent to $n-1$ mod $2$. 
\end{proposition}
\begin{proof}
    Observe from the tables that every $v_{1}$-torsion free generator in
    \[
    \Ext_{\cE(1)_{\star}}(L(\nu_{2}(k!)), L(\nu_{2}(m!))
    \]
    satisfies the statement. Furthermore, the statement is preserved under action by an element in $\Ext_{\cE^{\mR}_\star(1)}(\mM_{2}, \mM_{2})$ (see \cref{tab:ExtM2M2gens}). Finally, suspension by the $C_2$-regular representation also preserves the statement. So indeed every $v_{1}$-torsion free element of 
    \[
    E_{2}^{s,f,w} \cong \bigoplus\limits_{k=0}^{\infty}\Ext_{\cE(1)_{\star}}(L(\nu_{2}(k!)), L(\nu_{2}(m!))
    \]
    has a motivic weight fitting the desired description.
\end{proof}

\subsubsection{Differentials from $v_{1}$-torsion free classes to $v_{1}$-torsion classes}
We will now show there are no differentials with source a $v_1$-torsion free class and target a $v_{1}$-torsion class. The $v_1$-torsion classes in filtration $f > 0$ are precisely the classes in the $v_0$-towers of the triangle formation (\cref{lem:k>m} (1)) with $f >0$. In particular, these classes occur only in odd stem $s \leq -3$. Since 
\[
\bigoplus\limits_{k \le m}\Sigma^{\rho(m-k)}\Ext_{\cE(1)}( L(\nu_{2}(k!) ), L(\nu_{2}(m!) ) )
\]
is contained entirely in stem $s \ge 0$, it only remains to show that there are no differentials originating in $v_1$-torsion free classes of even stem in 
\[
\bigoplus\limits_{k > m}\Sigma^{\rho(m-k)}\Ext_{\cE(1)}( L(\nu_{2}(k!) ), L(\nu_{2}(m!) ) )
.\]


Fix $n \ge 0$. We will put bounds on the Adams filtration and motivic weight of generators in stem $-2n$. Then we will put bounds on the Adams filtration and motivic weight of potential $v_1$-torsion targets, which are in stem $-2n-1$. This will allow us to rule out the potential for any nonzero Adams differentials. 

The only even-stem $v_1$-torsion free generator in the positive cone summand of 
\[
\Sigma^{\rho(m-k)}\Ext_{\cE(1)}( L(\nu_{2}(k!) ), L(\nu_{2}(m!) ) )
\]
is $\Sigma^{\rho(m - k)} x$, which has stem $2(m-k)$ and filtration $\nu_{2}(k!) - \nu_{2}(m!)$ (see \cref{tab:Extgens k>m pos}). Similarly, the only even-stem $v_1$-torsion free generator in the negative cone summand is $\Sigma^{\rho(m - k)} x'$ (see \cref{tab:Extgens k>m neg}), which has stem $2(m-k)$ and filtration $\nu_{2}(k!) - \nu_{2}(m!) - 1$. So in \[
\bigoplus\limits_{k > m}\Sigma^{\rho(m-k)}\Ext_{\cE(1)}( L(\nu_{2}(k!) ), L(\nu_{2}(m'!) ) ),\]
there are exactly two $v_{1}$-torsion free generators in degree $-2n$: one positive cone generator $\Sigma^{-\rho n}x$, which has filtration $\nu_{2}(k!) - \nu_{2}(k-n)!$, and one negative cone generator $\Sigma^{-\rho n}x'$, which has filtration $\nu_{2}(k!) - \nu_{2}(k-n)!-1$.

One can check in the same manner that at stem $-2n-1$, any $v_{1}$-torsion element in the positive cone component of 
\[
\bigoplus\limits_{m>k}\Sigma^{(m-k)\rho}\Ext_{\cE(1)_{\star}}\left(L(\nu_{2}(k!)), L(\nu_{2}(m!)) \right)
\]
has Adams filtration at most $\nu_{2}(k!) - \nu_{2}\left((k-n+1)! \right)$. Meanwhile, the negative cone summands are shifted down by $1$ in filtration: any $v_{1}$-torsion element in the negative cone component has Adams filtration at most $\nu_{2}(k!) - \nu_{2}\left((k-n+1)! \right)-1$. 

\begin{comment}
\textcolor{blue}{Proving these bounds proceeds similarly to the nonequivariant argument given in \cref{nonequiv: collapse}. Specifically, observe first that in $ \Ext_{\cE(1)_{\star}}\left(L(\nu_{2}(k!), L(\nu_{2}(m!)\right)$, the highest filtration of any $v_{1}$-torsion free class in stem $-3-2i$ is $\nu_{2}(k!)-\nu_{2}(m!)-i$ in the positive cone component. (This is recorded in \cref{lem:k>m}, but it is helpful to compare with \cref{fig:clExtL3Fp}). So in $\Sigma^{\rho(k-m)}\Ext_{\cE(1)_{\star}}\left(L(\nu_{2}(k!), L(\nu_{2}(m!)\right)$, the highest filtration of any $v_{1}$-torsion free class in stem $-3-2i - 2(k-m)$ is $\nu_{2}(k!)-\nu_{2}(m!)-i$ in the positive cone component. So for any $0 \le i \le m < k$ such that 
\[-3 -2i -2(k-m) = -2n -1,\]
the highest filtration class will be $\nu_{2}(k!) - \nu_{2}\left((k-n+1+i)!\right) - i$ in the positive cone. So at any stem $-2n-1$, the highest filtration class will be in filtration $\nu_{2}(k!) - \nu_{2}((k-n+1)!)$ in the positive cone. The upper bound on the Adams filtration of $v_1$-torsion classes in the negative cone component follows from a similar argument. }
\end{comment}


These bounds on Adams filtration rule out the possibility of any remaining nontrivial differentials, except possibly from the negative cone summand to the positive cone summand. However, we also rule these out for degree reasons, specifically by considering motivic weight.

Any generator in stem $-2n$ of the negative cone summand has motivic weight at least $-n + 4$, while any $v_{1}$-torsion class in stem $-2n-1$ of the the positive cone summand with filtration $f > 0$ has  motivic weight at most $-n$. 

\begin{comment}
This can be seen as follows: $\Sigma^{\rho(m-k)}\frac{\gamma}{\tau^{4j}}x'$ has motivic weight $-n+4(j+1)$. So the minimum motivic weight of any negative cone component generator in stem $-2n$ is indeed $-n+4$. Now consider a class $\Sigma^{\rho(m-k)}v_{0}^{i}y_{j}$ in $\Sigma^{\rho(k-m)}\Ext_{\cE(1)_{\star}}\left(L(\nu_{2}(k!), L(\nu_{2}(m!)\right)$. The stem of this class is $2(\nu_{2}(m!)-\nu_{2}(k!) + j + m-k)-1$, and the motivic weight of this class is $\nu_{2}(m!)-\nu_{2}(k!) + j + m -k$. So if a class $\Sigma^{\rho(m-k)}v_{0}^{i}y_{j}$ has stem $-2n-1$, then the motivic weight that class is $-n$. Thus the motivic weight of $\Sigma^{\rho(m-k)}\tau^{2r}v_{0}^{i}y_{j}$ is at most $-n$. 
So indeed no differentials from $v_{1}$-torsion free to $v_{1}$-torsion classes are possible. 
\end{comment}
\end{proof}

\end{document}
